\section{Parameters for a polynomial filter}

\begin{theorem}[A joint count and contraction estimate]
    \label{thm:chain_polynomial_filter_parameters}
    \lean{Real.exists_dg_scalar_parameters}
    \leanok
    For every $C_H>0$ and $M\in\R$, there is a real constant
    $A\ge\max\{1,21,M,\log400000\}$ with the following properties.
    For every integer $d\ge2$ and $0<g\le1$, put
    \begin{align}
        u    & =1+\log d+\log(1/g),                   &
        \ell & =\left\lceil Au^2/g\right\rceil,\notag                                           \\
        c_A  & =6+2\log(A+1),                         & \mathcal D & =(d\ell^2)^{22\ell}.\notag
    \end{align}
    Then $Au^2/g\le\ell\le(A+1)u^2/g$, $A\le\ell$,
    $20\ell+2\lfloor\ell^2/\ell\rfloor=22\ell$, and
    \begin{align}
        \log\mathcal D\le22(A+1)c_Ag^{-1}u^3.\notag
    \end{align}
    For every $0<R\le2C_H\ell$, define
    $\Delta=4\exp(-4\ell^2\sqrt{(4g/5)/R})$. One has
    $\mathcal D^{12}\Delta\le1/100000$.
    The same $A$ has the following simultaneous tail property. For all
    $C,c,\varepsilon>0$ and $p\ge0$ satisfying
    $\max\{0,\log(C/\varepsilon)\}+p+c\le cM$, one has
    $C\exp(-c(\ell-g))\le\varepsilon g^p$.
    This is a supporting scalar estimate for the parameter choice
    in~\cite[lines 1010--1024]{Arad2016RigorousRG}.
\end{theorem}
\begin{proof}\leanok
    Since $\log(A+1)=o(\sqrt A)$, one may choose $A$ above all the prescribed
    lower bounds so that
    \begin{align}
        264\bigl(6+2\log(A+1)\bigr)+1
         & \le4\sqrt{\frac{2A}{5C_H}}.\notag
    \end{align}
    The ceiling bounds follow from $u\ge1$ and $g\le1$. They imply
    $\log\ell\le\log(A+1)+2\log u+\log(1/g)$. Since $\log u\le u$,
    one obtains $\log(d\ell^2)\le c_Au$ and
    $12\log \mathcal D\le264\ell c_Au$.
    The upper ceiling bound gives the asserted bound on $\log\mathcal D$.
    Separately, since $\ell>0$, integer division gives
    $\lfloor\ell^2/\ell\rfloor=\ell$, proving the stated exponent identity.

    From $g\ell\ge Au^2$ and $R\le2C_H\ell$, comparison of nonnegative
    squares gives
    $u\sqrt{2A/(5C_H)}\le\ell\sqrt{(4g/5)/R}$. Consequently
    \begin{align}
        12\log \mathcal D-4\ell^2\sqrt{\frac{4g/5}{R}}
                              & \le-\ell u\le-A,\notag                            \\
        \mathcal D^{12}\Delta & \le4e^{-A}\le\frac4{400000}=\frac1{100000}.\notag
    \end{align}
    Finally put $B=\max\{0,\log(C/\varepsilon)\}$. The threshold condition,
    $\ell\ge Au$, and $u\ge1+\log(1/g)$ give
    \begin{align}
        c(\ell-g) & \ge cAu-c\ge(B+p+c)u-c\notag                             \\
                  & \ge B+p\log(1/g)\ge\log(C/\varepsilon)+p\log(1/g).\notag
    \end{align}
    Exponentiation proves the tail estimate with the same $A$.
\end{proof}

\begin{theorem}[One parameter choice for contraction and the physical tail]
    \label{thm:physical_dg_scalar_parameters}
    \lean{Real.exists_physical_dg_scalar_parameters}
    Let
    \begin{align}
        C_* & =288e^{11/8},\notag                              \\
        M_* & =8\bigl(\max\{0,\log(C_*/10^{-6})\}+1+1/8\bigr).
    \end{align}
    There exists a real $A\ge\max\{1,21,M_*,\log400000\}$ such that,
    for every integer $d\ge2$ and $0<g\le1$, the quantities
    \begin{align}
        u          & =1+\log d+\log(1/g), & \ell & =\lceil Au^2/g\rceil,\notag \\
        \mathcal D & =(d\ell^2)^{22\ell}, & c_A  & =6+2\log(A+1)
    \end{align}
    satisfy $21\le\ell$, $2\le\mathcal D$, $A\le\ell$,
    $Au^2/g\le\ell\le(A+1)u^2/g$, and
    \begin{align}
        \log\mathcal D     & \le22(A+1)c_Ag^{-1}u^3,\notag \\
        C_*e^{-(\ell-g)/8} & \le10^{-6}g.
    \end{align}
    For every $0<R\le8\ell$,
    \begin{align}
        \mathcal D^{12}\,4e^{-4\ell^2\sqrt{(4g/5)/R}}\le10^{-5}.
    \end{align}
    The same $A$ gives both estimates for all $d,g,R$.
\end{theorem}
\begin{proof}
    \uses{thm:chain_polynomial_filter_parameters}
    Apply the joint scalar parameter theorem with boundary constant $4$
    and lower threshold $M_*$. Its simultaneous tail estimate applies with
    $C=C_*$, $c=1/8$, error $10^{-6}$ and power $1$ by the definition of
    $M_*$. The ceiling lower bound gives $\ell\ge21$ and $\mathcal D\ge2$.
    This is the scalar choice used for the physical specialization of
    \cite[Theorem t:ais, lines 1010--1024]{Arad2016RigorousRG}.
\end{proof}

\begin{theorem}[Endpoints for the five-region Chebyshev interval]
    \label{thm:physical_chebyshev_endpoints}
    \lean{Real.physical_chebyshev_endpoints}
    For $\ell\ge21$, $0\le E_0\le4$ and $0<g\le1$, put
    $\eta_0=E_0+g/10$, $\eta_1=E_0+9g/10$ and $R=7\ell-\eta_0$.
    Then $\eta_0<\eta_1\le7\ell$, $0<R\le8\ell$,
    $\eta_1-\eta_0=4g/5$, and
    \begin{align}
        \bigl(2e^{-2\ell^2\sqrt{(\eta_1-\eta_0)/R}}\bigr)^2
        =4e^{-4\ell^2\sqrt{(4g/5)/R}}.
    \end{align}
\end{theorem}
\begin{proof}
    Substitute the definitions. The bounds on $E_0,g,\ell$ make the interval
    nonempty and give the upper bound for $R$. Square the exponential.
\end{proof}

\begin{theorem}[Absorbing a collar at square polynomial degree]
    \label{thm:physical_dg_collar_budget}
    \lean{TNLean.MPS.AreaLaw.collar_card_le_square_at_square_degree}
    Let $q\ge2$, $\ell\ge1$, $d\le2\ell$ and $n$ be natural numbers.
    If $n\le((q\ell^2)^{18\ell+2\lfloor\ell^2/\ell\rfloor})^2$, then
    \begin{align}
        nq^{2d}\le((q\ell^2)^{22\ell})^2.
    \end{align}
\end{theorem}
\begin{proof}
    Since $\ell>0$, the original exponent is $20\ell$.
    The collar contributes $q^{2d}\le(q\ell^2)^{4\ell}$, which is absorbed
    by the difference between $40\ell$ and $44\ell$ in the squared bound.
\end{proof}

\begin{theorem}[One scalar choice for all accuracy parameters]
    \label{thm:dg_accuracy_scalar_parameters}
    \lean{Real.exists_dg_accuracy_scalar_parameters}
    Put $C_*=288e^{11/8}$, $C_+=C_*e^{1/8}$ and
    $M_+=8(\max\{0,\log(C_+/10^{-6})\}+1+1/8)$.
    There is a real constant $A\ge\max\{1,21,M_+,\log400000\}$ such that,
    for every integer $m\ge2$ and $0<g\le1$, the quantities
    \begin{align}
        u_m          & =1+\log m+\log(1/g),     & \ell_m & =\lceil Au_m^2/g\rceil,\notag \\
        \mathcal D_m & =(m\ell_m^2)^{22\ell_m}, & c_A    & =6+2\log(A+1)
    \end{align}
    satisfy $21\le\ell_m$, $m\le\mathcal D_m$, $2\le\mathcal D_m$,
    $A\le\ell_m$, $Au_m^2/g\le\ell_m\le(A+1)u_m^2/g$, and
    \begin{align}
        \log\mathcal D_m     & \le22(A+1)c_Ag^{-1}u_m^3,\notag \\
        C_*e^{-(\ell_m-g)/8} & \le\frac{g}{10^6m}.
    \end{align}
    For every $0<R\le8\ell_m$,
    $\mathcal D_m^{12}\,4e^{-4\ell_m^2\sqrt{(4g/5)/R}}\le10^{-5}$.
    One also has $20\ell_m+2\lfloor\ell_m^2/\ell_m\rfloor=22\ell_m$.
    The same $A$ works for all $m,g,R$.
\end{theorem}
\begin{proof}
    \uses{thm:chain_polynomial_filter_parameters}
    Apply the joint scalar theorem with boundary constant $4$ and threshold
    $M_+$. To obtain the accuracy factor, put $g'=g/m$. Then
    $1+\log(1/g')=u_m$, $Au_m\le\ell_m$ and $0<g'\le1$.
    The scalar tail estimate with coefficient $C_+$ gives
    $C_+e^{-(\ell_m-g')/8}\le10^{-6}g'$.
    Since $g-g'\le1$, the desired tail with coefficient $C_*$ is no larger.
    The ceiling and logarithmic estimates are those of the joint scalar
    choice. The separate integer identity follows from $\ell_m>0$.
    This is a numerical refinement for the physical interval families;
    it does not change the on-site Hilbert space.
\end{proof}

\begin{theorem}[A collar count with a larger numerical base]
    \label{thm:collar_larger_count_base}
    \lean{TNLean.MPS.AreaLaw.collar_card_le_square_of_larger_count_base}
    Let $2\le q\le m$, $\ell\ge1$, $d\le2\ell$ and $n$ be natural numbers.
    If $n\le((q\ell^2)^{18\ell+2\lfloor\ell^2/\ell\rfloor})^2$, then
    \begin{align}
        nq^{2d}\le((m\ell^2)^{22\ell})^2.
    \end{align}
    Only the numerical bound uses $m$; the physical collar still has $q$
    states per site.
\end{theorem}
\begin{proof}
    \uses{thm:physical_dg_collar_budget}
    First absorb the collar using the bound with base $q$.
    Then increase this numerical base from $q$ to $m$ in its nonnegative
    integer power.
\end{proof}
